\documentclass{ximera}

\input{../../preamble.tex}

\title{Theory}

\begin{document}

\begin{abstract}
the exponential story
\end{abstract}
\maketitle


Exponential functions exhibit a constant percentage growth rate. The same change anywhere in the domain produces the same percengatge change in the function value. That is the function value at the beginning of the domain interval is multiplied by the same constant to produce the function value at the end of the domain interval.

\textbf{Core Exponential Functions}

We have two types of core exponential functions, $Exp(x) = r^x$. 


Their different behaviors are dictated by the base.



$\blacktriangleright$ If $r > 1$, then $Exp(x) = r^x$ is an increasing exponential function.

$\blacktriangleright$ If $0 < r < 1$, then $Exp(x) = r^x$ is a decreasing exponential function.






The Elementary Functions include \textbf{general exponential functions}, which insert a leading coefficient.






\begin{definition} \textbf{\textcolor{green!50!black}{Exponential Functions}}

The \textbf{general exponential function} is a function that \textbf{\textcolor{purple!85!blue}{can}} be represented in the form 

\[  Exp(x) = A \cdot r^x    \, | \,  A \ne 0, \, r > 0 \, \text{ with } \, r \ne 1   \]


We refer to \textbf{\textcolor{purple!85!blue}{A}} as the leading coefficient.

We refer to \textbf{\textcolor{purple!85!blue}{r}} as the base.


\end{definition}





The leading coefficient controls vertical stretching or compression. The sign of the leading coefficient dictates the sign of our function values, which is illustrated graphically as a vertical reflection. 






\begin{image}
\begin{tikzpicture}
   \begin{axis}[name = leftgraph, 
            domain=-10:10, ymax=10, xmax=10, ymin=-10, xmin=-10,
            axis lines =center, xlabel=$x$, ylabel={$y$},
            ticklabel style={font=\scriptsize},
            every axis y label/.style={at=(current axis.above origin),anchor=south},
            every axis x label/.style={at=(current axis.right of origin),anchor=west},
            axis on top
          ]
          
          \addplot [line width=1, gray, dashed,samples=200,domain=(-10:10),<->] {0};


          \addplot [line width=2, penColor, smooth,samples=100,domain=(-9:5.7)] {2 * (1.3^x)};
          \addplot [line width=2, penColor, smooth,samples=100,domain=(-9:4.4)] {3 * (1.3^x)};
          \addplot [line width=2, penColor, smooth,samples=100,domain=(-9:2.4)] {5 * (1.3^x)};
          \addplot [line width=2, penColor, smooth,samples=100,domain=(-9:8)] {0.3 * (1.3^x)};
          \addplot [line width=2, penColor2, smooth,samples=100,domain=(-9:8), <->] {1.3^x};

          \addplot [color=penColor2,only marks,mark=*] coordinates{(0,1)};

          \addplot [line width=2, penColor, smooth,samples=100,domain=(-9:5.7)] {-2 * (1.3^x)};
          \addplot [line width=2, penColor, smooth,samples=100,domain=(-9:4.4)] {-3 * (1.3^x)};
          \addplot [line width=2, penColor, smooth,samples=100,domain=(-9:2.4)] {-5 * (1.3^x)};
          \addplot [line width=2, penColor, smooth,samples=100,domain=(-9:8)] {-0.3 * (1.3^x)};
          \addplot [line width=2, penColor2, smooth,samples=100,domain=(-9:8), <->] {-1 * 1.3^x};

          \addplot [color=penColor2,only marks,mark=*] coordinates{(0,-1)};




  \end{axis}
\end{tikzpicture}
\end{image}






The base dictates a growing or decaying function - the constant prcentage change could be an increase or decrease.






\begin{image}
\begin{tikzpicture}
   \begin{axis}[name = leftgraph, 
            domain=-10:10, ymax=10, xmax=10, ymin=-10, xmin=-10,
            axis lines =center, xlabel=$x$, ylabel={$y$},
            ticklabel style={font=\scriptsize},
            every axis y label/.style={at=(current axis.above origin),anchor=south},
            every axis x label/.style={at=(current axis.right of origin),anchor=west},
            axis on top
          ]
          
          \addplot [line width=1, gray, dashed,samples=200,domain=(-10:10),<->] {0};


          \addplot [line width=2, penColor, smooth,samples=100,domain=(-5.7:9)] {2 * (1.3^(-x))};
          \addplot [line width=2, penColor, smooth,samples=100,domain=(-4.4:9)] {3 * (1.3^(-x))};
          \addplot [line width=2, penColor, smooth,samples=100,domain=(-2.4:9)] {5 * (1.3^(-x))};
          \addplot [line width=2, penColor, smooth,samples=100,domain=(-8:9)] {0.3 * (1.3^(-x))};
          \addplot [line width=2, penColor2, smooth,samples=100,domain=(-8:9), <->] {1.3^(-x)};

          \addplot [color=penColor2,only marks,mark=*] coordinates{(0,1)};

          \addplot [line width=2, penColor, smooth,samples=100,domain=(-5.7:9)] {-2 * (1.3^(-x))};
          \addplot [line width=2, penColor, smooth,samples=100,domain=(-4.4:9)] {-3 * (1.3^(-x))};
          \addplot [line width=2, penColor, smooth,samples=100,domain=(-2.4:9)] {-5 * (1.3^(-x))};
          \addplot [line width=2, penColor, smooth,samples=100,domain=(-8:9)] {-0.3 * (1.3^(-x))};
          \addplot [line width=2, penColor2, smooth,samples=100,domain=(-8:9), <->] {-1 * 1.3^(-x)};

          \addplot [color=penColor2,only marks,mark=*] coordinates{(0,-1)};




  \end{axis}
\end{tikzpicture}
\end{image}

















Graphs of exponential functions have the horizontal axis as a horizontal asymptote in one direction or the other, which illustrates that they have no zeros.

\begin{explanation} \textbf{\textcolor{red!75!green}{Base greater than $1$}} 


When $r > 1$ we have a \wordChoice{\choice{decaying} \choice[correct]{growing}} function. It could be growing positively or negatively, depending on the sign of theleading coefficient.

\[  \lim_{x \to -\infty} A \, r^x = 0 \, \text{ and } \, \lim_{x \to \infty} A \, r^x = \pm\infty \]


Since $r^x$ is always positive, it is the sign of $A$ that dictates if the unbounded growth is positive or negative.




\begin{image}
\begin{tikzpicture}
   \begin{axis}[name = leftgraph, 
            domain=-10:10, ymax=10, xmax=10, ymin=-10, xmin=-10,
            axis lines =center, xlabel=$x$, ylabel={$A>0$},
            every axis y label/.style={at=(current axis.above origin),anchor=south},
            every axis x label/.style={at=(current axis.right of origin),anchor=west},
            axis on top
          ]
          
          \addplot [line width=1, gray, dashed,samples=200,domain=(-10:10),<->] {0};

          \addplot [line width=2, penColor, smooth,samples=100,domain=(-9:8), <->] {1.3^x};
          \addplot [color=penColor,only marks,mark=*] coordinates{(0,1)};

           

  \end{axis}
  \begin{axis}[at={(leftgraph.outer east)},anchor=outer west, 
            domain=-10:10, ymax=10, xmax=10, ymin=-10, xmin=-10,
            axis lines =center, xlabel=$x$, ylabel={$A<0$},
            every axis y label/.style={at=(current axis.above origin),anchor=south},
            every axis x label/.style={at=(current axis.right of origin),anchor=west},
            axis on top
          ]
          
          \addplot [line width=1, gray, dashed,samples=200,domain=(-10:10),<->] {0};

          \addplot [line width=2, penColor, smooth,samples=100,domain=(-9:8),<->] {-(1.3^x)};
          \addplot [color=penColor,only marks,mark=*] coordinates{(0,-1)};

           

  \end{axis}
\end{tikzpicture}
\end{image}


Graphs of exponential functions have one strategic point, which occurs when the exponent equals zero. For the basic exponential function, this point is $\left( \answer{0}, A \right)$.


When the base is greater than $1$, exponential functions have unbounded growth in the direction that makes the exponent large and \wordChoice{\choice[correct]{positive} \choice{negative}}.


When the base is greater than $1$, exponential functions approach $0$ in the direction that makes the exponent large and \wordChoice{\choice{positive} \choice[correct]{negative}}.




\end{explanation}





\begin{explanation} \textbf{\textcolor{red!75!green}{Base less than $1$}}



When $r < 1$, we have a decaying function. The horizontal axis is still a horizontal asymptote, just in the direction where the exponent is large and positive. The function now decays. It could decay positively or negatively, depending on the sign of leading coefficient.


\[ \lim_{x \to -\infty} A \, r^x = \pm\infty \, \text{ and } \, \lim_{x \to \infty} A \, r^x = 0 \]


 
The sign of $A$ dictates if the function decays through positive or negative values, because $r^x > 0$.



\begin{image}
\begin{tikzpicture}
   \begin{axis}[name = leftgraph, 
            domain=-10:10, ymax=10, xmax=10, ymin=-10, xmin=-10,
            axis lines =center, xlabel=$x$, ylabel={$A>0$},
            every axis y label/.style={at=(current axis.above origin),anchor=south},
            every axis x label/.style={at=(current axis.right of origin),anchor=west},
            axis on top
          ]
          
          \addplot [line width=1, gray, dashed,samples=200,domain=(-10:10),<->] {0};

          \addplot [line width=2, penColor, smooth,samples=200,domain=(-8:9), <->] {1.3^(-x)};
          \addplot [color=penColor,only marks,mark=*] coordinates{(0,1)};
           

  \end{axis}
  \begin{axis}[at={(leftgraph.outer east)},anchor=outer west, 
            domain=-10:10, ymax=10, xmax=10, ymin=-10, xmin=-10,
            axis lines =center, xlabel=$x$, ylabel={$A<0$},
            every axis y label/.style={at=(current axis.above origin),anchor=south},
            every axis x label/.style={at=(current axis.right of origin),anchor=west},
            axis on top
          ]
          
          \addplot [line width=1, gray, dashed,samples=200,domain=(-10:10),<->] {0};

          \addplot [line width=2, penColor, smooth,samples=200,domain=(-8:9),<->] {-(1.3^(-x)};
          \addplot [color=penColor,only marks,mark=*] coordinates{(0,-1)};
           

  \end{axis}
\end{tikzpicture}
\end{image}



Graphs of basic exponential functions have one strategic point, which occurs when the exponent equals zero. For the basic exponential function, this point is $(0, A)$, which is below the horizontal asymptote, when $A < 0$.




When the base is less than $1$, exponential functions have unbounded growth in the direction that makes the exponent large and \wordChoice{\choice[correct]{positive} \choice{negative}}.


When the base is less than $1$, exponential functions approach $0$ in the direction that makes the exponent large and \wordChoice{\choice{positive} \choice[correct]{negative}}.


\end{explanation}






All exponential functions share a common structure.


\begin{itemize}
\item Each is always positive or always negative. This depends on the sign of the leading coefficient.
\item Each approaches $0$ in one direction.
\item In the other direction, they grow unbounded (postively or negatively). 
\end{itemize}





All exponential graphs share a common structure.


\begin{itemize}
\item Each is always above or below the horizontal axis. This depends on the sign of the leading coefficient.
\item Each has the horizontal axis as an asymptote in one direction.
\item In the other direction, they extend unbounded (up or down). 
\end{itemize}




These are the important aspects or characteristics that we watch when moving onto shifting and stretching the graphs of exponential functions.












\begin{summary} \textbf{\textcolor{red!70!yellow}{Exponential Behavior}}


Basic exponential functions, $E(x) = A \cdot r^x$, are either increasing functions or decreasing functions.


$\blacktriangleright$  \textbf{\textcolor{purple!85!blue}{Base Greater than $1$: $r > 1$}} 


\begin{itemize}
\item greater positive exponents mean multiplying by the base more, which results in larger values.  
\item greater negative exponents mean multiplying by the reciprocal of the base more, which results in smaller values.  
\end{itemize}


The leading coefficient in front, $A$, tells us if this larger/smaller value is positive or negative.


\begin{itemize}
\item $A > 0$ and $r > 1$ : \wordChoice{\choice[correct]{increasing} \choice{decreasing}}  (positive) function
\item $A < 0$ and $r > 1$ : \wordChoice{\choice{increasing} \choice[correct]{decreasing}}  (negative) function  
\end{itemize}






$\blacktriangleright$  \textbf{\textcolor{purple!85!blue}{Base Less than $1$: $0 < r < 1$}}  


\begin{itemize} 
\item greater positive exponents mean multiplying by the base more, which results in smaller values.  
\item greater negative exponents mean multiplying by the reciprocal of the base more, which results in larger values.  
\end{itemize}


The leading coefficient in front, $A$, tells us if this larger/smaller value is positive or negative.


\begin{itemize}
\item $A > 0$ and $r < 1$ : \wordChoice{\choice{increasing} \choice[correct]{decreasing}}  (positive) function
\item $A < 0$ and $r < 1$ : \wordChoice{\choice[correct]{increasing} \choice{decreasing}}  (negative) function  
\end{itemize}



\end{summary}


























\subsection*{Analysis}

As we see, basic exponential functions can be

\begin{itemize}
\item positive and increasing ($A>0$ and $r>1$)
\item positive and decreasing ($A>0$ and $r<1$)
\item negative and increasing ($A<0$ and $r<1$)
\item negative and decreasing ($A<0$ and $r>1$)
\end{itemize}






















\begin{example}  Analysis



Analyze   $f(x) = \frac{1}{3} \cdot \, 2^{x+5}$ 




\textbf{\textcolor{red!75!green}{explanation}} 




Categorize: $f$ is an exponential function, because it can match our template for exponential functions, $A \, r^x$.


\[
\frac{1}{3} \cdot \, 2^{x+5} = \frac{1}{3} \cdot \, 2^{x} \cdot 2^{5} = \frac{32}{3} \cdot \, 2^{x}
\]


\textbf{Just Thinking...}  

We'll begin by thinking about the graph to help guide our analysis. 






The ``inside'' of the formula for $f(x)$ is $x+5$.  This equals $0$, when $x=-5$.  The exponent is positive for $x>-5$, and because the base is $2 > 1$, this is the direction (right) of unbounded growth.  Therefore, the other direction (left) is where the horizontal asymptote is in effect.  Since the leading coefficient is $\frac{1}{3} > 0$, the unbounded growth is positive.

At $x=-5$, we have our one anchor point for the graph.  The point is $\left(-5, \frac{1}{3} \right)$, which is $\frac{1}{3}$ above the horizontal axis.


Graph of $y = f(x)$.

\begin{image}
\begin{tikzpicture}
  \begin{axis}[
            domain=-10:10, ymax=10, xmax=10, ymin=-10, xmin=-10,
            axis lines =center, xlabel=$x$, ylabel=$y$, grid = major,
            ytick={-10,-8,-6,-4,-2,2,4,6,8,10},
            xtick={-10,-8,-6,-4,-2,2,4,6,8,10},
            ticklabel style={font=\scriptsize},
            every axis y label/.style={at=(current axis.above origin),anchor=south},
            every axis x label/.style={at=(current axis.right of origin),anchor=west},
            axis on top
          ]
          
          \addplot [line width=1, gray, dashed,samples=200,domain=(-10:10),<->] {0};

           \addplot [line width=2, penColor, smooth,samples=200,domain=(-10:0.6),<->] {0.33 * 2^(x+5)};

           \addplot[color=penColor,fill=penColor,only marks,mark=*] coordinates{(-5,0.333)};

          


 

  \end{axis}
\end{tikzpicture}
\end{image}








\textbf{\textcolor{red!75!green}{Now, for our (algebraic) function analysis.}} 


\textbf{\textcolor{purple!65!blue}{Algebraic Analysis}}


\textbf{Domain:}  The domain of every exponential function is $(-\infty, \infty)$. 




\textbf{Zeros:}  Exponential functions do not have zeros. 




\textbf{Continuity:}  Exponential functions are continuous.  



\textbf{End-Behavior:} 

The base of $f$ is $2 > 1$. 



$2^{positive}$ will get big. $2^{negative}$ will get small. 

The linear exponent of $f$ is $x+5$ has a positive leading coefficient. 

The negative direction is the direction that makes the exponent large and negative.  That makes this the direction that the function value approaches $0$. 


\[
\lim\limits_{x \to -\infty} f(x) = 0
\]




The other direction is where the exponential function is uunbounded.  Since the leading coefficient of $f$ is positive, $f$ is unbounded positively. 



\[
\lim\limits_{x \to \infty} f(x) = \infty
\]









\textbf{Behavior:} 

We can think of $f$ as a composition with two linear functions.


\begin{itemize}
\item Let $L_{out}(u) = \frac{1}{3} \, u$, an increasing linear function, since it is a linear function with a positive leading coefficient.
\item Let $E(t) = 2^t$, an increasing core exponential function, since it is a core exponential function with a base greater than $1$.
\item Let $L_{in}(v) = v + 5$, an increasing linear function, since it is a linear function with a positive leading coefficient.
\end{itemize}





\[  f(x) = \frac{1}{3} \cdot \, 2^{x+5} = (L_{out} \circ E \circ L_{in})(x)\]

\[ inc \circ inc \circ inc = inc  \]



This story gets shortened to leading coefficients and the base:\\
The base of $f$ is $2$, which is greater than $1$. The leading coefficient of $f$ is $\frac{1}{3}$, which is positive. The leading coefficient of the linear exponent is $1$, which is also positive.  These make $f$ is a positive increasing exponential function.



\textbf{Global Extrema:}  Exponential functions do not have global maximums or minimums. 


\textbf{Local Extrema:}  Exponentials functions do not have local maximums or minimums. 



\textbf{Range:} The end-behavior tells us that the range is $(0, \infty)$.





\end{example}

















$\blacktriangleright$ \textbf{\textcolor{blue!55!black}{Shifting or Not}} 



Exponential functions can swallow up horizontal graphical shifting into the leading coefficient.



If we add/subtract $b$ in the exponent, for a horizontal stretch, then we have

\[
a \cdot r^{x + b} = a \cdot r^x \cdot r^b = a \cdot r^b \cdot r^x = (a \cdot r^b) \cdot r^x
\]

We obtain a new leading coefficient. 

It just depends on how you like to see exponential functions. 

















$\blacktriangleright$ \textbf{\textcolor{blue!55!black}{Stretching or Not}} 



Exponential functions swallow up graphical stretching and compressing into the leading coefficient and base.


\begin{itemize}
\item If we multiply on the outside, by $b$, for a vertical stretch, then we have

\[
b \cdot (a \cdot r^x) = (a \cdot b) \cdot r^x
\]

We obtain a new leading coefficient. 






\item If we multiply on the inside, by $b$, for a horizontal stretch, then our exponential rules give us

\[
a \cdot r^{b \cdot x} = a \cdot (r^b)^x = a \cdot R^x
\]

We obtain a new base, which tells us about the growth of the function.


\end{itemize}


It just depends on how you like to see exponential functions. 







Algebra let's us see exponential functions in more complex forms.








\begin{definition} \textbf{\textcolor{green!50!black}{More General Exponential Functions}}

An \textbf{exponential function} is any function that \textbf{\textcolor{purple!85!blue}{CAN}} be represented with a formula of the form

\[     Exp(x) =    A \cdot r^{B \, x + C}          \]

where $A$, $B$, and $C$ are real numbers and $A \ne 0$, $B \ne 0$, $r > 0$ with $r \ne 1$.


\end{definition}


\begin{center}
\textbf{\textcolor{red!70!black}{We can use any algebra, any time, any where that we feel helps us.}}  
\end{center}



\begin{idea}  \textbf{\textcolor{blue!55!black}{Exponential Compositions}} 

We can view exponential functions as compositions


\begin{itemize}
\item Let $L_{out}(u) = A \, u $, a linear function
\item Let $E(t) = r^t$, a core exponential function
\item Let $L_{in}(v) = B \, v + C$, a linear function
\end{itemize}

With these core component functions, we can view exponential functions as compositions.


\[  Exp(x) =    A \cdot r^{B \, x + C} = (L_{out} \circ E \circ L_{in})(x) \]


Now the Chain Rule can give us the behavior.

\end{idea}






































\begin{example} Analysis



Analyze $G(t) = 4 \cdot 3^{2t-1}$.   


\textbf{\textcolor{red!75!green}{explanation}} 



Categorize:  $G$ is an exponential function, since it matches our more general official template $A \, r^{B \, x + C}$. 


\textbf{Just Thinking...}

We could take advantage of exponential algebra to simplify the exponent.   

\[ 
G(t) = 4 \cdot 3^{2t-1} = 4 \cdot 3^{2t} \cdot \frac{1}{3} = \frac{4}{3} \cdot (3^2)^t  = \frac{4}{3} \cdot 9^t
\]


We now have a simpler formula with the new base $\answer{9}$ and the new leading coefficient is $\answer{\frac{4}{3}}$.




The leading coefficient is $\frac{4}{3} > 0$ and the base is $9 > 1$. Therefore, $G(t)$ is a positive increasing exponential function.



$G(t) = \frac{4}{3} \cdot 9^t - 5$ is a transformation of $g(t) = 3^t$.  Let's compare $G(t)$ back to $3^t$.


Since the new base $9$ is greater than the old base of $3$, the function grows faster.  The graph goes up to the right faster and steeper.  This means the graph is compressed horizontally, which is the result of multiplying by $2$ in the exponent.



\textbf{\textcolor{red!75!green}{With these ideas, we can create an algebraic analysis.}} 


\textbf{\textcolor{purple!65!blue}{Algebraic Analysis}}



\textbf{Domain:}  $G$ is an exponential function, so its domain is $(-\infty, \infty)$.





\textbf{Zeros:}  $G$ is an exponential function, so it has no zeros.




\textbf{Continuity:}  $G$ is an exponential function, so it is continuous.




\textbf{Behavior}


$\blacktriangleright$ Thinking about the formula pieces.

$G(t) = 4 \cdot 3^{2t-1}$ is an exponential function.  It  is either always increasing or always decreasing.


The leading coeffcient is $4$, which is positive.  Therefore, $G$ is positive.

The leading coefficient of the linear exponent is $2$, which is positive.  That with a base of $3 > 1$ tell us us that $G$ becomes unbounded towards the positive side of the real numbers.


That means $G$ is an increasing function.



$\blacktriangleright$ Using the Chain Rule


\begin{itemize}
\item Let $L_{out}(u) = 4 \, u$, an increasing linear function, since the leading coefficient is positive.
\item Let $E(x) = 3^x$, an increasing core exponential function, since the base is greater than $1$.
\item Let $L_{in}(v) = 2 \, v - 1$, an increasing linear function, since the leading coefficient is positive.
\end{itemize}

With these component functions, we can view $G$ as a composition.


\[  G(t) = 4 \cdot 3^{2t-1} = (L_{out} \circ E \circ L_{in})(t) \]

\[ inc \circ inc \circ inc = inc  \]


$G$ is an increasing function.



\textbf{End-Behavior}


$G$ is a positive increasing function.



\[
\lim\limits_{t \to -\infty}G(t) = \answer{0}  
\]

\[
\lim\limits_{t \to \infty}G(t) = \answer{\infty}
\]





\textbf{Local Maximum and Minimum:}  Exponential functions do not have local maximums or minimums.






\textbf{Global Maximum and Minimum:}  Exponential functions do not have global maximums or minimums.




\textbf{Range:}  $G$ is a positive exponential function.   Its range is $(0, \infty)$.



\end{example}





\begin{warning}  \textbf{\textcolor{red!70!black}{Sign vs. Behavior}}


Exponential functions are either always positive or always negative. The sign of the leading coeffcient dictates this.


However, behavior is not so simple.


Behavior is controlled by three pieces of the formula: the outside linear function, the core exponential, and the inside linear function.


\begin{itemize}
   \item the leading coefficient
   \item the base
   \item the leading coefficient of the exponent (a linear function)
\end{itemize}



\textbf{For Instance,} $f(x) = 3 \cdot e^{x}$ is an increasing function.

\textbf{\textcolor{blue!55!black}{$\blacktriangleright$}} If we change the sign of the leading coefficient, $g(x) = -3 \cdot e^{x}$ is a decreasing function.


\textbf{\textcolor{blue!55!black}{$\blacktriangleright$}} Or, if instead, we change the sign of the exponent's leading coefficient, $h(x) = 3 \cdot e^{-x}$ is a decreasing function.


\textbf{\textcolor{blue!55!black}{$\blacktriangleright$}} Or, if instead, we change the base to be less than $1$, $h(x) = 3 \cdot \left( \frac{1}{2} \right)^{x}$ is a decreasing function.



We must watch all three to determine the behavior.


This is exactly what the Chain Rule tells us.

\end{warning}




\begin{observation} \textbf{\textcolor{-red!75}{Big Bases : $r > 1$}}     



When multiplying by $B$ on the inside of an exponential function, the graph experiences a horizontal stretch or compression.


If $1 < B$, then we are compressing horizontally.  This works out in the algebra. When $1 < B$, then $r < r^B$. The base becomes larger, which means the function grows faster, which means the graph goes up faster, which means the graph has been compressed horizontally. 



If $0 < B < 1$, then we are stretching horizontally.  This works out in the algebra. When $0 < B < 1$, then $r^B < r$. The base becomes smaller, which means the function grows slower, which means the graph goes up slower, which means the graph has been stretched horizontally.


\end{observation}





One way to watch this is to think of how the exponential function is built from a core exponential function through composition.


Begin with the base, is it greater than or less than $1$?  That will tell you if the core exponential function is increasing or decreasing.

Then look at the signs of the two leading coefficients. If they have the same sign, then the behavior doesn't change. If they have different signs, then the behavior changes.


















\begin{example}  Analysis



Analyze   $B(t) = -2 \, \left( \frac{2}{3} \right)^{3-t}$ 



\textbf{\textcolor{red!75!green}{explanation}} 



Categorize: $B$ is an exponential function, since it matches our template, $A \, r^{B \, x + C} $. 





\textbf{Just Thinking...}  

$\blacktriangleright$  the base: $\frac{2}{3} < 1$   \\
$\blacktriangleright$  the leading coefficient is negative.  \\
$\blacktriangleright$  the exponent: $3-t$ is positive for large negative $t$. \\
$\blacktriangleright$  the exponent: $3-t$ is negative for large positive $t$. 



Together, these tell us that $B(t)$ settles down for large negative values of $t$ and that $B(t)$ becomes unbounded when $t$ is large and positive.



This is a transformed version of the core exponential function template $\left( \frac{2}{3} \right)^{-t} = \left( \answer{\frac{3}{2}} \right)^t$.  



When $t < 0$, then $-t > 0$ and we get  $\left( \frac{2}{3} \right)^{positive}$ and the core exponential portion is becoming smaller, approaching $0$.  





\[ \lim\limits_{t \to -\infty} \left( \frac{2}{3} \right)^{-t}  = 0 \]



When $t > 0$, then $-t < 0$ and we get  $\left( \frac{2}{3} \right)^{negative} = \left( \frac{3}{2} \right)^{positive}$ and $B(t)$ is becoming unbounded.  



\[ \lim\limits_{t \to \infty} \left( \frac{2}{3} \right)^{-t}  = \lim\limits_{t \to \infty} \left( \frac{3}{2} \right)^t  = \infty \]




Exponential growth to the right and decay to the left.






Since the leading coefficient is $-2 < 0$, these smaller/larger values of $-2 \, \left( \frac{2}{3} \right)^{-t}$ are smaller/larger negative values.






Horizontal Shift

Our exponent here is $3 - t$.  Our function's exponent is zero, $3-t=0$ when $t=3$. Our one anchor point is shifted over to $3$.  Multipying by $-2$, means the dot is $2$ away(below) from the horizontal asymptote, which is $y=0$.  Our anchor point is $(3, -2)$.








Graph of $y = B(t)$.

\begin{image}
\begin{tikzpicture}
  \begin{axis}[
            domain=-10:10, ymax=10, xmax=10, ymin=-10, xmin=-10,
            axis lines =center, xlabel=$t$, ylabel=$y$, grid = major,
            ytick={-10,-8,-6,-4,-2,2,4,6,8,10},
            xtick={-10,-8,-6,-4,-2,2,4,6,8,10},
            ticklabel style={font=\scriptsize},
            every axis y label/.style={at=(current axis.above origin),anchor=south},
            every axis x label/.style={at=(current axis.right of origin),anchor=west},
            axis on top
          ]
          
          \addplot [line width=1, gray, dashed,samples=200,domain=(-10:10),<->] {0};

           \addplot [line width=2, penColor, smooth,samples=200,domain=(-10:7.5),<->] {-2 * (0.666^(3-x))};

           \addplot[color=penColor,fill=penColor,only marks,mark=*] coordinates{(3,-2)};

          


  \end{axis}
\end{tikzpicture}
\end{image}





\textbf{\textcolor{red!75!green}{With these ideas guiding us, we can write an algebraic analysis. }} 





\textbf{\textcolor{blue!55!black}{Algebraic Analysis}} 






\textbf{Domain:}  $B$ is an exponential function, so its domain is $(-\infty, \infty)$.





\textbf{Zeros:}  $B$ is an exponential function, so it has no zeros.




\textbf{Contiuity:}  $B$ is an exponential function, so it is continuous.




\textbf{Behavior}



$\blacktriangleright$ \textbf{Thinking with the Chain Rule...}


\begin{itemize}
\item Let $L_{out}(u) = -2 \, u$, a decreasing linear function, since the leading coefficient is negative.
\item Let $E(x) = \left( \frac{2}{3} \right)^x$, a decreasing core exponential function, since the base is less than $1$. 
\item Let $L_{in}(v) = 3 - v$, a decreasing linear function, since the leading coefficient is negative.
\end{itemize}


\[
B(t) = -2 \, \left( \frac{2}{3} \right)^{3-t} = (L_{out} \circ E \circ L_{in})(t)
\]

\[ dec \circ dec \circ dec = dec \]




$\blacktriangleright$ \textbf{Thinking of the formula pieces...}


$B(t) = -2 \, \left( \frac{2}{3} \right)^{3-t}$ is an exponential function.  It  is either always increasing or always decreasing.


The leading coeffcient is $-2$, which is negative.  Therefore, $B$ is negative.

The leading coefficient of the linear exponent is $-1$, which is negative.  That with a base of $\frac{2}{3} < 1$ tell us us that $B$ becomes unbounded towards the positive side of the real numbers.


$B$ is a negative function that becomes unbounded toward the positive end of the real numbers.  That means $G$ is a decreasing function.



Or, we can compare back to a core exponential function...


$B$ is an exponential function, which might be compared to the core exponential function $e^x$, which is a positive increasing function.


First, the base has been changed from $e > 1$ to $\frac{2}{3} < 1$.  That switches the behavior to decreasing.

Second, the leading coefficient of the linear exponent is $-1$, which is negative.  That switches the behavior again, back to increasing.


Third, the leading coefficient is $-2 < 0$. That switches the behavior again, back to decreasing.







\textbf{End-Behavior}


$B$ is a negative decreasing function.



\[
\lim\limits_{t \to -\infty}B(t) = \answer{0}  
\]

\[
\lim\limits_{t \to \infty}B(t) = \answer{\infty}
\]





\textbf{Local Maximum and Minimum:}  Exponential functions do not have local maximums or minimums.






\textbf{Global Maximum and Minimum:}  Exponential functions do not have global maximums or minimums.




\textbf{Range:}  $B$ is a negative exponential function.   Its range is $(-\infty, 0)$.







\end{example}

























In the example above, we could have algebraically moved the horizontal shift to the leading coefficient.



\[
-2 \, \left( \frac{2}{3} \right)^{3-t}
\]


\[
-2 \, \left( \frac{2}{3} \right)^{3} \cdot \left( \frac{2}{3} \right)^{-t}
\]


\[
-2 \, \left( \frac{8}{27} \right)  \cdot \left( \frac{2}{3} \right)^{-t}
\]


\[
\left( \frac{-16}{27} \right)  \cdot \left( \frac{2}{3} \right)^{-t}
\]


We could have further transfered the negative sign in the exponent to the base by reciprocating the base.

\[
\left( \frac{-16}{27} \right)  \cdot \left( \frac{3}{2} \right)^t
\]



Algebra provides many tools for modifying the representing formula and altering how we think about the behavior. \\


It just depends on how you like to think of it. \\





\begin{example}  Graphs



Analyze   $K(f) = 3^{5-f}$ 


\begin{question}
 

The base is \wordChoice{\choice{less than}\choice[correct]{greater than}} $1$.
\end{question}
\begin{question}



\begin{question}

The linear exponent is \answer{$5-f$}

\end{question}


The exponent gets big and positive when $f$ gets big and \wordChoice{\choice{positive}\choice[correct]{negative}}.
\end{question}
\begin{question}


The graph will become unbounded to the \wordChoice{\choice{right}\choice[correct]{left}}.\\
\end{question}

\begin{question}


The horizontal asymptote is $y = \answer{0}$.
\end{question}

\begin{question}


The graph will approach the asymptote to the \wordChoice{\choice[correct]{right}\choice{left}}.\\
\end{question}
\begin{question}


Our one strategic anchor point moves to $\left(\answer{5}, \answer{1}\right)$.
\end{question}
\begin{question}


The graph will become unbounded \wordChoice{\choice[correct]{up}\choice{down}}.\\
\end{question}




Graph of $y = K(f)$.

\begin{image}
\begin{tikzpicture}
  \begin{axis}[
            domain=-10:10, ymax=10, xmax=10, ymin=-10, xmin=-10,
            axis lines =center, xlabel=$f$, ylabel=$y$, grid = major,
            ytick={-10,-8,-6,-4,-2,2,4,6,8,10},
            xtick={-10,-8,-6,-4,-2,2,4,6,8,10},
            ticklabel style={font=\scriptsize},
            every axis y label/.style={at=(current axis.above origin),anchor=south},
            every axis x label/.style={at=(current axis.right of origin),anchor=west},
            axis on top
          ]
          
          \addplot [line width=1, gray, dashed,samples=200,domain=(-10:10),<->] {0};

           \addplot [line width=2, penColor, smooth,samples=200,domain=(2.75:10),<->] {3^(5-x)};

           \addplot[color=penColor,fill=penColor,only marks,mark=*] coordinates{(5,1)};

 

  \end{axis}
\end{tikzpicture}
\end{image}





\end{example}





















\begin{center}
\textbf{\textcolor{green!50!black}{ooooo-=-=-=-ooOoo-=-=-=-ooooo}} \\

more examples can be found by following this link\\ \link[More Examples of Exponential Functions]{https://ximera.osu.edu/csccmathematics/precalculus2/precalculus2/expFunctions/examples/exampleList}

\end{center}









\end{document}
